\chapter{Detailed Considerations, Challenges, and Pitfalls in Accurate Memory Measurement}
\label{ch:appendix-challenges-pitfalls}

This appendix provides an extensive discussion of the complexities, pitfalls, and potential sources of error that arise when attempting to measure memory consumption accurately in Python applications.
It highlights both internal and external factors that can distort memory metrics, and it explains the methodologies used to mitigate these issues.
It also describes various experiments, including the generation of synthetic seismic datasets, the isolation of execution environments through Docker, the use of diverse memory measurement tools, and the evaluation of memory pressure constraints.


\section{Introduction}
\label{sec:appendix-introduction}

Accurate memory measurement in Python is often challenging, especially in scientific computing and \ac{HPC} contexts.
Resource allocation in these environments directly affects performance, scalability, and cost-effectiveness.
A clear understanding of an application’s memory consumption informs decisions on how to best allocate hardware resources and can reveal inefficiencies in data handling or algorithmic structure.
However, several aspects of Python’s runtime behavior, as well as external operating system optimizations, introduce complexity into the measurement process.

Python manages memory dynamically.
Its interpreter frequently allocates and deallocates memory through reference counting, augmented by cyclical garbage collection.
These mechanisms do not necessarily free memory immediately upon object deletion.
This can lead to transient peaks in usage that inflate measured consumption.
In addition, libraries such as NumPy or TensorFlow can allocate memory outside of Python’s garbage collector, making purely Python-centric approaches insufficient to capture the total footprint.
Operating systems also introduce complexities like \ac{COW}, memory compression (e.g., zram~\cite{zram}), and caching optimizations that obscure application-level usage.

When memory measurements are gathered in shared or interactive environments (e.g., Jupyter notebooks), previously allocated objects might linger unexpectedly in memory from prior executions.
These artifacts can bias results unless the environment is carefully reset or fully isolated between runs.
Furthermore, measurement tools themselves add overhead and can subtly alter the memory usage patterns they aim to measure.
Thus, both the methodology for data collection and the environment in which the application runs must be controlled to ensure valid experimental outcomes.


\section{Common Challenges and Pitfalls}
\label{sec:appendix-common-challenges}

\subsection{Dynamic Memory Allocation and Fragmentation}
\label{subsec:appendix-dynamic-allocation}

Python’s dynamic allocation routines and internal memory manager can introduce significant variability in memory consumption.
Garbage collection occurs periodically and does not necessarily release memory back to the operating system as soon as objects become unreferenced.
Additionally, Python’s memory allocator (\texttt{pymalloc}) may cause internal fragmentation, resulting in unoccupied regions of memory that remain allocated but unavailable for reuse by other processes.

Fragmentation is a particular concern when large structures are allocated and deallocated multiple times in unpredictable sequences.
In \ac{HPC} and data analysis scenarios, where massive arrays (often in the order of gigabytes) can be created and destroyed frequently, fragmentation can accumulate in ways that inflate measured memory usage.

\subsection{Delayed Garbage Collection}
\label{subsec:appendix-delayed-gc}

In many experiments, residual memory can persist from previous workloads if a Python process remains active.
This is often observed in Jupyter notebooks, which allow repeated execution of cells without reloading the interpreter.
Garbage collection itself might be triggered at non-deterministic intervals or based on thresholds related to object allocations and deallocations.
Thus, code that deallocates memory at a logical level may not immediately translate into reduced system-level memory usage.

When measuring memory consumption in a repeated-run workflow, it is essential to account for these delayed deallocations.
Otherwise, memory metrics from later runs may reflect additional usage from earlier operations.
This kind of carryover can be particularly misleading if experiments are run in quick succession, causing spurious or inflated peaks.

\subsection{Operating System and Container-Level Effects}
\label{subsec:appendix-os-effects}

Modern operating systems include virtual memory systems that abstract physical RAM into a continuous address space for processes.
Features like \ac{COW} can postpone actual memory usage until writes occur.
Linux systems also offer zswap~\cite{zswap} or zram~\cite{zram} to compress memory pages, which complicates direct measurements of usage at the process level.

Containers introduce another layer of abstraction.
Although Docker and other containerization technologies isolate processes effectively, the memory usage reported inside a container can be slightly different from that observed at the host level.
Containers can cap available memory and trigger \ac{OOM} termination if usage exceeds certain thresholds.
Nevertheless, containers also rely on the host kernel’s memory management subsystem, so container-level metrics must sometimes be cross-referenced with host-level data to gain a full understanding.

\subsection{Caching Mechanisms and Reused Objects}
\label{subsec:appendix-caching}

At both the Python interpreter and operating system level, various caches can persist across multiple runs of a program.
For example, Python maintains internal caches of small integers, certain strings, and recently deallocated memory blocks.
The operating system may cache disk blocks in \ac{RAM}.
These caches are beneficial for performance, but they interfere with straightforward memory measurement because they introduce unaccounted usage that fluctuates unpredictably over time.

Third-party libraries, particularly large-scale numeric frameworks, also maintain internal buffers and caches.
These caches might not be visible to Python’s built-in measurement methods and might require specialized profiling at the native layer (e.g., hooking into C/C++ allocations).

\subsection{Measurement Tool Overheads}
\label{subsec:appendix-measurement-tool-overheads}

Tools that profile memory usage can themselves introduce overhead.
For instance, sampling-based profilers increase \ac{CPU} load by periodically querying memory states.
Tools that track individual allocations and deallocations may store their own metadata in memory.
This overhead can inflate measured memory usage or alter the timing of application execution, thus changing garbage collection behavior.

When possible, it is advisable to verify results with multiple measurement strategies.
Comparisons between lightweight external monitors (e.g., Docker’s own metrics) and Python-internal profilers (e.g., \texttt{tracemalloc}) can detect significant divergences that might be due to the profiler itself.


\section{Overview of Experimentation Notebooks}
\label{sec:appendix-experimentation-overview}

A series of Jupyter notebooks were created to investigate these challenges in practice~\footnote{Experiment Folder. Available at \url{https://github.com/discovery-unicamp/memory-aware-chunking/blob/main/experiments/00-appendix-a-pitfalls-and-limitations-of-memory-profiling-on-linux/notebooks}. Accessed: 2025-05-06.}.
Each notebook targets a specific stage of the experimental workflow: data generation, environment isolation, multi-method memory profiling, and controlled memory pressure scenarios.
Although they were executed sequentially, each can stand alone to illustrate distinct issues related to accurate memory measurement.

\subsection{Notebook 1: Data Generation}
\label{subsec:appendix-notebook-data-generation}

This notebook focuses on generating synthetic seismic datasets.
A synthetic seismic volume is created in \ac{SEG-Y} format with dimensions such as \mbox{$100 \times 100 \times 100$}, ensuring consistency across experiments.
The shape is logged, and inline slices of the volume are displayed as grayscale heatmaps to validate the coherence of the synthetic data.

%\begin{figure}[h]
%    \centering
%    \includegraphics[width=0.6\textwidth]{placeholder-data-generation}
%    \caption{Example placeholder showing a grayscale heatmap of one inline section from a synthetic 3D seismic dataset.}
%    \label{fig:appendix-data-generation}
%\end{figure}

As seen in Figure~\ref{fig:appendix-data-generation}, darker shades correspond to lower amplitudes and lighter shades to higher amplitudes.
Although the data is synthetically generated, the layered patterns resemble real-world seismic reflectors.
This notebook confirms that data generation proceeds reliably and reproducibly.
No external interferences or prior memory usage lingered here since each dataset was created afresh in a controlled environment.

\subsection{Notebook 2: Execution Environment and Isolation}
\label{subsec:appendix-notebook-execution-environment}

The second notebook examines how Docker containers can be leveraged to isolate experiments and eliminate memory contamination between consecutive runs.
Repeated Docker executions of a moderately sized seismic operation were performed—often an envelope extraction routine.

Memory usage in each run was logged, showing peaks ranging from roughly \mbox{4.9 GB} to \mbox{5.8 GB} across multiple iterations.
Because each test was executed in a fresh container, memory states from previous runs did not carry over.
Repeated runs demonstrated consistent peak measurements, confirming that container-based isolation effectively prevents memory cross-contamination.

Although the notebook primarily presents numeric logs, an imagined plot of memory usage versus iteration count would show a relatively narrow distribution of peak values across 30 or more runs, demonstrating reproducibility.

%\begin{figure}[h]
%    \centering
%    \includegraphics[width=0.7\textwidth]{placeholder-container-iterations}
%    \caption{Placeholder illustration of peak memory usage across repeated containerized runs.}
%    \label{fig:appendix-container-iterations}
%\end{figure}

\subsection{Notebook 3: Memory Profiling Across Tools}
\label{subsec:appendix-notebook-memory-profiling}

The third notebook focuses on comparing multiple measurement tools:
\begin{itemize}
    \item \textbf{psutil} (process-level data including RSS and VMS),
    \item \textbf{resource} (peak memory tracking),
    \item \textbf{tracemalloc} (object-level tracing inside Python),
    \item \textbf{/proc filesystem} queries (OS-level insights).
\end{itemize}

Each approach was run in parallel or in separate trials to validate consistency.
Results showed that psutil’s reported memory usage closely aligned with Docker metrics for overall container usage, while \texttt{tracemalloc} offered granular visibility into Python-specific allocations.
Discrepancies arose in cases where large native libraries (e.g., NumPy) performed allocations outside Python’s control.
In such instances, the \texttt{/proc} or Docker-based measurements captured broader usage than \texttt{tracemalloc} did.

%\begin{figure}[h]
%    \centering
%    \includegraphics[width=0.65\textwidth]{placeholder-tool-comparison}
%    \caption{Placeholder figure comparing memory usage data from different profiling tools over time.}
%    \label{fig:appendix-tool-comparison}
%\end{figure}

When correlated with external data from Docker’s API, minor overheads introduced by high-frequency sampling in \texttt{tracemalloc} were noted.
This overhead did not exceed a few percent of total memory usage in these tests, but it could be more significant in memory-constrained scenarios or extremely large-scale HPC jobs.

\subsection{Notebook 4: Memory Pressure Scenarios}
\label{subsec:appendix-notebook-memory-pressure}

The final notebook applied controlled memory constraints to the same envelope extraction algorithm (and other seismic routines) to validate measured memory peaks.
Initially, a baseline run without memory caps revealed a peak usage of about \mbox{5.33 GB}.
Subsequent runs introduced progressive limitations: first near 5.0 GB, then 4.9 GB, and so forth.

Under constraints that were slightly below the observed peak, failures occurred.
In other words, runs were terminated by Docker’s OOM killer or signaled by Python exceptions.
These outcomes confirmed that a roughly \mbox{5.3 GB} measurement was not merely an artifact, but a true lower bound on the memory needed for successful operation under the tested conditions.

%\begin{figure}[h]
%    \centering
%    \includegraphics[width=0.6\textwidth]{placeholder-memory-pressure}
%    \caption{Placeholder figure showing how the measured peak memory usage corresponds to actual memory limits that cause task failure.}
%    \label{fig:appendix-memory-pressure}
%\end{figure}

By plotting memory usage over time under increasingly strict memory caps, the point of failure could be identified with high precision.
This type of memory constraint validation supports the reliability of the initial measurements.


\section{Selected Algorithms and Their Memory Profiles}
\label{sec:appendix-selected-algorithms}

\subsection{Envelope}
\label{subsec:appendix-env}

The Envelope algorithm computes the instantaneous amplitude of a seismic trace by obtaining the magnitude of the analytic signal.
Its computational demands are relatively low and do not typically induce large memory burdens, making it a clear baseline test.
In memory measurement exercises, Envelope provided a straightforward way to confirm that measurement tools produce stable and consistent results under light workloads.

\subsection{\ac{GST3D}}
\label{subsec:appendix-gst3d}

The \ac{GST3D} algorithm applies time-frequency analysis to three-dimensional seismic data.
It is typically associated with heavy calculations and sizable memory allocations, since large 3D arrays must be transformed and processed repeatedly.
Testing \ac{GST3D} revealed how memory usage can climb sharply when dealing with large seismic volumes, exposing fragmentation and garbage collection issues more readily than lightweight tasks do.

\subsection{3D Gaussian Filtering}
\label{subsec:appendix-gaussian}

Three-dimensional Gaussian filtering is commonly used in signal and image processing for smoothing or noise reduction.
It involves convolving a 3D matrix with a Gaussian kernel, which can be memory intensive for large volumes.
Since it is not strictly a seismic algorithm, it served as a general-purpose workload for evaluating memory measurement techniques in a different domain.
Results from 3D Gaussian filtering were consistent with those from \ac{GST3D} in terms of revealing peaks and intermediate memory usage fluctuations.


\section{Experimental Isolation and Validation Techniques}
\label{sec:appendix-isolation-validation}

\subsection{Importance of Container Isolation}
\label{subsec:appendix-container-isolation}

All experiments were executed within Docker containers to ensure that memory usage from one experiment could not leak into the next.
Isolation mitigated the possibility of residual allocations or delayed garbage collections from a previous test interfering with subsequent measurements.
Each container was launched from a clean state, with dedicated CPU and memory allocations.
This approach proved vital for generating reproducible and unbiased memory metrics, especially in repeated-run scenarios.

\subsection{Cross-Verification with Docker API}
\label{subsec:appendix-cross-verification}

Each internal measurement approach (psutil, resource, tracemalloc, \texttt{/proc} inspections) was cross-validated with memory usage reported by Docker’s own metrics.
This external vantage point allowed for the detection of measurement bias introduced by instrumentation overhead or sampling frequency.
In general, alignment was observed within a few percentage points, but \texttt{tracemalloc} consistently reported lower values than Docker for total process usage, since tracemalloc focuses primarily on Python-level object allocations rather than system-level memory footprints.

\subsection{Failure-Based Validation of Peak Usage}
\label{subsec:appendix-failure-validation}

Additional confirmation of measured peak usage was obtained by imposing memory caps just below previously recorded peaks.
When the process failed under these constraints, it corroborated that the measured peak was genuine.
A memory-constrained environment ensured that measurement anomalies, such as artificially high reported peaks that could be ignored by the application, were less likely.


\section{Summary of Observed Issues and Recommendations}
\label{sec:appendix-issues-recommendations}

\subsection{Residual Memory from Previous Runs}
\label{subsec:appendix-residual-memory}

A principal source of erroneous data was memory leftover from prior runs.
Even explicitly deleting Python objects did not guarantee immediate release of memory to the operating system.
To address this pitfall:
\begin{itemize}
    \item Restarting the Python interpreter or Jupyter kernel before each measurement is recommended.
    \item Alternatively, isolating each test in a fresh Docker container is highly effective.
\end{itemize}

\subsection{Misleading Metrics from Single-Tool Measurements}
\label{subsec:appendix-misleading-metrics}

Relying on only one measurement technique can be misleading.
Each tool yields different perspectives:
\begin{itemize}
    \item Tools like \texttt{resource} capture only peak RSS, which is useful but does not reveal how usage evolves over time.
    \item Python-level tracers like \texttt{tracemalloc} do not account for memory allocated outside the interpreter.
    \item External system monitors (e.g., Docker, \texttt{/proc}) provide a holistic view but lack Python-specific granularity.
\end{itemize}

When computational tasks are large-scale or involve native libraries, cross-verification with multiple tools is strongly advised.

\subsection{Tool Overhead and Garbage Collection Timing}
\label{subsec:appendix-tool-overhead}

Some profiling methods alter the frequency of garbage collection or add their own overhead memory usage.
This overhead can slightly inflate measured metrics.
If repeated sampling over short intervals is used, overhead could become more substantial.
For memory-critical workloads, it may be necessary to reduce the sampling rate or rely on a combination of high-level external measures and carefully timed internal snapshots.

\subsection{Environment Pressure and Real-World Relevance}
\label{subsec:appendix-environment-pressure}

Peak usage measured in an unconstrained environment can differ from usage measured under memory pressure.
In actual HPC systems, resource constraints often induce different runtime behaviors, such as forced page swaps or more frequent garbage collections.
Imposing artificial memory caps during tests, as seen in the memory pressure experiments, helps mimic real-world conditions where machine resources are finite and shared among multiple jobs.


\section{Conclusion}
\label{sec:appendix-conclusion}

Accurate memory measurement for Python applications in HPC and scientific workflows demands an awareness of language-level abstractions, garbage collection, OS virtualization, and container-level optimizations.
Experimental protocols that involve container isolation, multiple measurement tools, cross-verification with external monitors, and memory constraint scenarios provide a robust framework for identifying genuine memory usage.
Nevertheless, pitfalls such as memory fragmentation, delayed garbage collection, and caching mechanisms necessitate careful planning of experimental design.
By adopting environment isolation (e.g., Docker containers), cross-checking with several profiling tools, and validating results under artificial memory caps, the risk of misinterpretation is substantially reduced.
It is recommended that future experiments incorporate all these layers of validation to ensure reliable insights into memory requirements and performance bottlenecks.